\scanpage{17}
\mspage{2}
(viewed as a component of such a ``generalised boundary''). We now should \deleted{should} get a modular space for ``compact conformal oriented surfaces with generalized boundary'' (of type $g,\nu$ and rigidified via $X_o$), call it $\widetilde{MB}_{g,\nu}$, homeomorphic to $\underline C^d\mathbin{\mathrm{x}}(\underline R^{\deleted{*}+})^{\nu}$, where now the second factor corresponds to \deleted{\unsure{the}} the ``parameters'' $\ell_i$, which are allowed to take also value $0$ (which means that the corresponding component of the generalized boundary is just one point). Thus $\widetilde{MB}_{g,\nu}$ appears as a variety with boundary (in the topological sense - in the real analytic sense, the ``boundary'' admits ``corner-like'' points obviously), and $\widetilde M_{g,\nu}$ appears as a part of the boundary.

My interest is in a better geometric understanding of the situation, which should be ``intrinsic'' namely not depend on any particular choice of a surgical decomposition of the reference surface $X_o$ into ``pants'', used in order to describe in a handy way \unsure{the} standard ``coordinate\deleted{\unsure{s}}functions'' on the modular space $\widetilde{MB}_{g,\nu}$. There appears to be a geometrically meaningful retraction of $\widetilde{MB}_{g,\nu}\deleted{\unsure{\nu}}$ upon $\widetilde M_{g,\nu}$ (commuting to the operations of the Teichmüller modular group), the fibers being homeomorphic to $(\underline R^+)^{\unsure{\nu}}$ - more specifically, I expect the semi-group $(\underline R^+)^I$ (where I is the set of indices for the ``holes'' of $X_o$) to act on $MB$ in a natural way, with free action of the subgroup $(R^{+*})^{\nu}$ upon $\widetilde{MB}^{o}$, in such a way that $\widetilde M$ is just the quotient of $\widetilde{MB}$ by this action (or of $\widetilde{MB}^{o}$ by the action of the \deleted{\unsure{unit-g}} corresponding subgroup), and that each fiber\inserted{$F$} is isomorphic to $(\underline R^+)^I$ by the choice of any ``origin'' in $F\cap\widetilde{MB}^{o}$.

Of course, ``\unsure{computionally}'', in terms of a decomposition of $X_o$ into pants, the idea of such an operation is pretty obvious - namely letting the components $\lambda_i$ of $\lambda\in(\underline R^+)^I$ act as a ``multiplier'' on the corresp\inserted{ond}ing coordinate $\ell_i$. However, it is not clear that this operation is intrinsic - and if it were intrinsic, an intrinsic geometric description would still be \deleted{\unsure{to}} desired.

Of course, in the description of the situation proposed above, the retraction of $\widetilde{MB}$ upon $\widetilde M$ is obtained by multiplying with the $0$ multiplier (all $\lambda_i$ are $0$). Now there is a direct geometrical description of \deleted{this} a retraction, by hyperbolic surgery. Namely, for any \inserted{compact} conformal surface of type $g,\nu$ with generalized boundary, let's ``fill in'' the holes which correspond to ordinary components of the boundary, which are riemannian oriented circles, by ``gluing in'' the cones on these circles (which are canonically endowed with a conformal structure, using the riemannian structure on the given circles). Thus we get a ``functor'' from \inserted{compact} conformal surfaces with generalized boundary (of type $g,\nu$) to compact conformal surfaces \uline{without boundary}, endowed with a system of $\nu$ points (making up a
% archive foliation: 16
